\section{Prim's Algorithm (Minimum Spanning Tree)}
\label{sec:graph-prim}

\problemstatement{Same as Section~\ref{sec:graph-kruskal}: find a
Minimum Spanning Tree of a weighted, connected, undirected graph --
this time by \textbf{growing a single tree} outward from one starting
node, rather than considering all edges globally sorted by weight.}

\begin{patternbox}
\emph{``Grow a tree, always extending it as cheaply as possible''}
mirrors the Graph Shortest Path chapter's Dijkstra's algorithm almost exactly: a
priority queue always offers the \textbf{cheapest available edge}
connecting some node already in the tree to some node not yet in it.
The key difference from Dijkstra: the priority queue tracks the
\textbf{edge weight to join the tree}, not the \textbf{cumulative
distance from a source} -- a node already close to the tree via one
cheap edge is preferred even if a different, longer path from the
start node exists.
\end{patternbox}

\subsection*{Why this also satisfies the cut property}

At every step, the tree built so far and the remaining nodes form a
cut (a partition into two groups). Prim's always picks the
\textbf{cheapest edge crossing that specific cut} -- exactly the
greedy choice the cut property (Section~\ref{sec:graph-kruskal})
guarantees is always safe to include in \emph{some} MST. Kruskal's
considers cuts implicitly through global edge sorting; Prim's
considers one specific, growing cut at each step.

\subsection*{Algorithm}

\begin{enumerate}
  \item Maintain a $\textit{inMST}$ boolean array and a min-heap of
        $(\textit{weight}, \textit{node})$ pairs, starting with
        $(0, \textit{anyStartNode})$.
  \item Pop the smallest-weight entry; if its node is already in the
        MST, skip (stale entry, exactly as in Dijkstra). Otherwise,
        mark it in the MST and add its weight to the total.
  \item For every neighbor of the newly added node not yet in the
        MST, push $(\textit{edgeWeight}, \textit{neighbor})$.
  \item Repeat until every node is in the MST.
\end{enumerate}

\subsection*{C++ implementation}

\begin{lstlisting}
#include <bits/stdc++.h>
using namespace std;

int primMST(int n, vector<vector<pair<int,int>>>& adj) {
    vector<bool> inMST(n, false);
    priority_queue<pair<int,int>, vector<pair<int,int>>, greater<>> pq;
    pq.push({0, 0}); // start from node 0, weight 0
    int totalWeight = 0;

    while (!pq.empty()) {
        auto [weight, node] = pq.top();
        pq.pop();
        if (inMST[node]) continue; // stale entry

        inMST[node] = true;
        totalWeight += weight;

        for (auto& [neighbor, edgeWeight] : adj[node]) {
            if (!inMST[neighbor]) {
                pq.push({edgeWeight, neighbor});
            }
        }
    }
    return totalWeight;
}

int main() {
    int n = 4;
    vector<vector<pair<int,int>>> adj(n);
    auto addEdge = [&](int u, int v, int w) {
        adj[u].push_back({v, w});
        adj[v].push_back({u, w});
    };
    addEdge(0, 1, 4); addEdge(0, 2, 1); addEdge(1, 2, 2);
    addEdge(1, 3, 5); addEdge(2, 3, 8);

    cout << primMST(n, adj) << endl; // 8
    return 0;
}
\end{lstlisting}

\subsection*{Dry run}

\dryrun{Same graph as Section~\ref{sec:graph-kruskal}'s example,
starting from node $0$.}

Heap: $[(0,0)]$. Pop $(0,0)$: add $0$ to MST, total $=0$; push
neighbors $(4,1)$ and $(1,2)$. Heap: $[(1,2),(4,1)]$. Pop $(1,2)$
(smallest): add $2$, total $=1$; push $2$'s neighbors not in MST:
$(2,1)$ (edge $1{-}2$, weight $2$) and $(8,3)$ (edge $2{-}3$, weight
$8$). Heap: $[(2,1),(4,1)_{\text{stale dup}},(8,3)]$. Pop $(2,1)$:
add $1$, total $=3$; push $1$'s remaining neighbor not in MST: $(5,3)$
(edge $1{-}3$, weight $5$). Heap now has both $(5,3)$ and $(8,3)$ for
node $3$. Pop $(4,1)_{\text{stale}}$: $1$ already in MST, skip. Pop
$(5,3)$ (smaller than the $(8,3)$ entry): add $3$, total $=8$. Final
MST weight: $8$ -- matching Kruskal's result on the same graph, as
expected (both algorithms always find \emph{a} minimum spanning tree,
and when the MST is unique, as here, the total weight necessarily
agrees).

\begin{complexitybox}
\textbf{Time Complexity.} $O(E \log V)$ -- each edge can trigger one
heap push, each costing $O(\log V)$.
\textbf{Space Complexity.} $O(V)$ for \textit{inMST}, plus $O(E)$ for
the heap in the worst case.
\end{complexitybox}

\begin{takeawaybox}
Kruskal's and Prim's solve the identical problem from two different
angles -- global edge sorting versus local frontier growth -- and
which is preferable depends on graph density: Kruskal's
$O(E\log E)$ favors sparse graphs (few edges to sort), while Prim's
$O(E\log V)$ with an adjacency-list-based heap favors graphs where $V$
is comparatively small relative to $E$. Both reuse machinery already
built earlier in this book: Kruskal's reuses Disjoint Set, Prim's
reuses Dijkstra's exact priority-queue frontier-expansion shape.
\end{takeawaybox}
